\subsection{A common finite word for the combined Hurwitz--Krammer tests}

\paragraph{D23.P14 (combined finite equations still admit a mixed state).}
There is an ordinary word $f\in F'''$ whose values simultaneously
satisfy the following finite conditions:
\begin{enumerate}
\item the complete Krammer pentagon of D23.P7, with first-pair
state $F_{0,1}$ and identity solvable blocks in the mixed pairs;
\item every two-cycle, three-cycle, reversed three-cycle and
$m=-1$ three-cycle in the full Nielsen closure of those Krammer
pairs, as checked in D23.C15 and D23.F24;
\item identity values on the entire Nielsen closure of all five
pairs in both Hurwitz representations of D23.P11;
\item the cubic detector value $\xi(f)=((1,1)^{\mathsf T},I)$
and the marked $S_4$ value $1$.
\end{enumerate}
Thus these combined finite equations admit a common-word mixed
state $(\rho,a_3,\nu_4)=(0,1,0)$.
The values also pass the mixed-pair generation test described
below, and the reflection conjugacy test in D23.C15's supplement.
No actual gentle automorphism or full-GT element is asserted.
In particular a simultaneous full-source lift of all peripheral
frames is not supplied by this proposition.

\begin{proof}
Let $S=\operatorname{SU}_3(3)$, the nonabelian simple group
of order $6048$. D23.P9 gives the exact image of $F'$ in the
first Krammer pair and the three marked mixed factors:
\[
 P_{12}'\times(H'\times S)^3.
\]
Here $|P_{12}|=46656$, $|H|=108$, and $H'''=1$.
Taking two further derived subgroups proves that the image
of $F'''$ in the three perfect coordinates is all of $S^3$,
while its image in every mixed $H$ coordinate is identity.
The same assertion holds for their full Nielsen closure,
because D23.C15 identifies exactly the same three perfect
kernel classes and D23.F24 identifies the corresponding
three solvable classes. The first-pair cyclic Nielsen values
are conjugates by $K_1K_2$, up to a central argument that
cancels on a commutator word.

Every evaluation-pair group in both Hurwitz representations
has order divisible only by $2$ and $3$. For the $648$-point
action this follows from the complete $66$-component order
inventory. For the $20736$-point action, the earlier inventory
computes all $372$ components of degree at most $216$;
\path{hurwitz_large_pair_orders.py} completes the remaining
$58$ exact orders. Every prime-to-six part is one.
No timed-out computation is included in this conclusion.
The three basic pair types suffice: reversal and the braid
half twist cover all five pentagon pairs, and a Nielsen
change of two generators does not change the group they
generate. Thus the full joint Hurwitz image, including all
these Nielsen markings, also has only the primes $2,3$.

Let $L$ be the joint image of $F'''$ in all these Hurwitz
coordinates, the first Krammer pair, all mixed solvable
coordinates, and the marked $S_4$ coordinate. Its order has
only the primes $2,3$. The common-source image in $L\times S^3$
projects onto both factors. There is no nontrivial common
quotient: every nontrivial quotient of $S^3$ is a product
of one or more copies of $S$, and hence has order divisible
by $7$, whereas every quotient of $L$ has order prime to $7$.
Goursat's lemma therefore gives the full product
\[
                 \operatorname{im}(F'''\to L\times S^3)
                    =L\times S^3.
\tag{D23.P14.1}
\]

The explicit word $f_*$ of D23.P12 supplies a particular
$L$-value: first Krammer state $F_{0,1}$, identity in every
Hurwitz Nielsen coordinate, identity in every mixed solvable
coordinate, and identity marked $S_4$ value. Among the P7
solutions labelled $(0,1)$ there are exactly three with
identity mixed solvable block. Choose any one of them and
take its three perfect coordinates. Equation (D23.P14.1)
gives one ordinary word $f\in F'''$ with this prescribed
$L$-value and these three perfect coordinates.

Its complete Krammer values are therefore exactly that P7
assignment. D23.C15 and D23.F24 verify all the stated Nielsen
cycle equations on every one of the $216$ assignments,
including these three. Its Hurwitz values are unchanged from
$f_*$ and so are identity on the entire finite closure.
The cubic value follows from the exact first-pair projection
of D23.P6 and the identification of $F_{0,1}$ in D23.F13;
the marked $S_4$ value was retained explicitly in $L$.
All assertions are consequently properties of one common
ordinary word, rather than independent assignments of the
same parameter in unrelated factors.
\end{proof}

\paragraph{A necessary generation test which needs no kernel rigidity.}
If $(1,f)$ is actual gentle, its normalized automorphism of
$F$ is $x\mapsto x$, $y\mapsto f^{-1}yf$. For every finite
surjection $q:F\twoheadrightarrow G=\langle X,Y\rangle$,
the composition $q\Phi_f$ is onto. Therefore necessarily
\[
                    \langle X,f(X,Y)^{-1}Yf(X,Y)\rangle=G.
\tag{D23.P14.2}
\]
This assertion does not require $\ker q$ to be invariant,
and is weaker than a claim that $\Phi_f$ induces an automorphism
of the particular quotient.

The exact program \path{krammer_gentle_generation.py} tests
(D23.P14.2) for both mixed Krammer pairs on all $216$ P7
assignments. In every case the solvable projection still has
order $108$ and the unitary projection still has order $24192$.
There are $100$ distinct block calculations after caching
identical finite matrices. The common abelian determinant
image is unchanged by conjugation, so the Goursat description
in D23.C10 shows that the entire original mixed pair group
is generated. Reversing a pair gives a conjugate generating
pair by the two-cycle, and the remaining Nielsen markings
are the already identified marked classes. All $36$ states
for each label $(\rho,a_3)$ pass this necessary test.

For the reflection supplement, the swap-stable perfect pair
has a literal intertwiner $J_0$ implementing $x\leftrightarrow y$,
with $J_0^2$ scalar. After projectivization it extends the
quotient to $F\rtimes\langle s\rangle$, $s^2=1$. The actual
identity $\beta s\beta^{-1}=sf$ consequently requires
$J_0f$ to be conjugate to $J_0$ by the pair image. The scalar
center of $U_3(3)$ does not affect conjugation, so an
$\operatorname{SU}_3(3)$ conjugator suffices. A projective
equality is also literal: comparing determinants forces a
scalar $\zeta$ with $\zeta^3=1$ in $\mathbb F_9^\times$, hence
$\zeta=1$. The checked conjugacy class has size $63$, and
all $216$ values pass. This is a necessary actual-source
frame test, not a GT-rigidity assertion for the Krammer kernel.

The result closes the finite equations and pointwise necessary
conditions listed above. It does not settle a global peripheral
frame lifting problem or decide A.84.
